\section{Теоретические основы алгоритмов сортировки}

\subsection{Общая классификация алгоритмов сортировки}

Алгоритмы сортировки делятся на устойчивые и неустойчивые, выполняющиеся на месте и требующие дополнительной памяти. В данной работе рассматриваются три представителя: пузырьковая сортировка (класс простых, $O(n^2)$), быстрая сортировка и сортировка слиянием (класс эффективных, $O(n \log n)$). Подробное описание этих алгоритмов приведено в фундаментальном труде~\cite{knuth1998}.

\subsection{Пузырьковая сортировка (Bubble Sort)}

Алгоритм многократно проходит по массиву, сравнивая соседние элементы и меняя их местами, если они расположены в неправильном порядке. После каждого прохода наибольший элемент «всплывает» в конец.

Число сравнений в худшем и среднем случаях:
\begin{equation}
    \label{eq:bubble_comp}
    C_{\text{bubble}} = \frac{n(n-1)}{2} = O(n^2)
\end{equation}

Число обменов также имеет квадратичную сложность. Несмотря на простоту, алгоритм неэффективен для больших $n$. Переменная \texttt{comparison\_count} обозначает количество выполненных сравнений.

\subsection{Быстрая сортировка (Quick Sort)}

Алгоритм «разделяй и властвуй»: выбирается опорный элемент (\texttt{pivot}), массив разбивается на элементы меньше опорного, равные ему и больше его, после чего рекурсивно сортируются подмассивы.

Средняя временная сложность:
\begin{equation}
    \label{eq:quick}
    T_{\text{quick}}(n) = O(n \log n)
\end{equation}
Детальный анализ быстрой сортировки представлен в работе~\cite{cormen2009}.

Худший случай (опорный элемент всегда минимум или максимум):
\begin{equation}
    \label{eq:quick_worst}
    T_{\text{quick}}^{\text{worst}}(n) = O(n^2)
\end{equation}

Быстрая сортировка выполняется на месте (в среднем требуется $O(\log n)$ дополнительной памяти для стека рекурсии).

\subsection{Сортировка слиянием (Merge Sort)}

Алгоритм также использует принцип «разделяй и властвуй»: массив рекурсивно делится пополам, каждая половина сортируется, затем происходит слияние двух упорядоченных массивов.

Гарантированная временная сложность:
\begin{equation}
    \label{eq:merge}
    T_{\text{merge}}(n) = O(n \log n)
\end{equation}
Теоретические основы сортировки слиянием также подробно рассмотрены в~\cite{cormen2009}.

В отличие от быстрой сортировки, сложность не зависит от входных данных. Недостаток "--- требуется дополнительная память $O(n)$ для хранения временных массивов при слиянии.

\subsection{Сравнительная таблица сложностей}

В таблице~\ref{tab:complexity} приведены сводные данные по трём алгоритмам.

\begin{table}[H]
\centering
\caption{Теоретическая сложность алгоритмов сортировки}
\label{tab:complexity}
\begin{tabular}{|l|c|c|c|c|}
\hline
\textbf{Алгоритм} & \textbf{Лучший случай} & \textbf{Средний} & \textbf{Худший} & \textbf{Память} \\
\hline
Пузырьковая & $O(n)$ & $O(n^2)$ & $O(n^2)$ & $O(1)$ \\
\hline
Быстрая & $O(n \log n)$ & $O(n \log n)$ & $O(n^2)$ & $O(\log n)$ \\
\hline
Слиянием & $O(n \log n)$ & $O(n \log n)$ & $O(n \log n)$ & $O(n)$ \\
\hline
\end{tabular}
\end{table}

\subsection{Ожидаемая разница между Python и C++}

Python является интерпретируемым языком с динамической типизацией, что накладывает дополнительные накладные расходы при каждой операции (проверка типов, управление памятью). C++ компилируется в машинный код, оптимизируется компилятором и работает непосредственно с аппаратными ресурсами. Ожидается, что C++ будет выполнять те же алгоритмы в \textbf{5--10 раз} быстрее, особенно при большом количестве операций сравнения и обмена.
